
\section{Architectural Assessment}

\subsection{Procedure Assessment}

\begin{center}
\begin{tabular}{|p{0.34\linewidth}|p{0.2\linewidth}|p{0.26\linewidth}|}
\hline
Procedure & Classification & Recommendation\\
\hline
\texttt{\_reduce\_hkl()} & Mathematical primitive & Keep; expand mathematical precondition in docstring.\\
\hline
\texttt{\_extended\_gcd()} & Mathematical primitive & Keep unchanged.\\
\hline
\texttt{\_right\_reduce\_row\_to\_hnf()} & Composite mathematical algorithm & Clarify docstring; evaluate decomposition into Bézout elimination primitives after audit.\\
\hline
\texttt{surface\_basis\_S\_from\_hkl()} & Orchestrator & Keep orchestration role explicit in comments.\\
\hline
\texttt{verify\_surface\_kernel\_is\_primitive()} & Mathematical primitive & Keep separate; document lattice-theoretic criterion.\\
\hline
\end{tabular}
\end{center}

\subsection{Documentation Assessment}

\begin{itemize}
\item State mathematical assumptions explicitly (primitive Miller index).
\item State mathematical guarantees separately from implementation details.
\item Explain orientation repair as gauge fixing.
\item Explain why each elementary transformation preserves unimodularity.
\end{itemize}

\subsection{Comments Assessment}

Comments should justify mathematical intent rather than restate code.
In particular, basis permutations and sign changes should be explained in
terms of lattice interpretation and gauge freedom.

\subsection{Refactoring Assessment}

No immediate refactoring is recommended before the audit is complete.

Potential future decomposition:

\begin{enumerate}
\item Primitive Miller reduction.
\item Bézout elimination step.
\item Unimodular completion.
\item Surface-basis interpretation.
\item Orientation gauge fixing.
\end{enumerate}

These recommendations are deferred until all transformation audits are
complete to avoid changing implementation while verification is in
progress.
